% ============================================================
%  preamble.tex
%  Packages, hyperref setup, theorem environments, project
%  macros and TikZ styles shared by the whole paper.
%  Included from paper.tex with % ============================================================
%  preamble.tex
%  Packages, hyperref setup, theorem environments, project
%  macros and TikZ styles shared by the whole paper.
%  Included from paper.tex with % ============================================================
%  preamble.tex
%  Packages, hyperref setup, theorem environments, project
%  macros and TikZ styles shared by the whole paper.
%  Included from paper.tex with % ============================================================
%  preamble.tex
%  Packages, hyperref setup, theorem environments, project
%  macros and TikZ styles shared by the whole paper.
%  Included from paper.tex with \input{preamble}.
% ============================================================

% ── Core packages ────────────────────────────────────────────
\usepackage[T1]{fontenc}
\usepackage[utf8]{inputenc}
\usepackage{microtype}          % better line breaking
\usepackage{cite}               % IEEE-style compressed citations [1]-[3]
\usepackage{amsmath,amssymb,amsthm}
\usepackage{graphicx}
\usepackage{booktabs}           % \toprule / \midrule / \bottomrule
\usepackage{calc}               % \widthof for fixed-width regime labels
\usepackage{algorithm}
\usepackage{algpseudocode}      % algorithmicx pseudocode
\usepackage{xcolor}
\usepackage{tikz}
\usetikzlibrary{positioning,arrows.meta,calc,shapes.geometric,fit,backgrounds}
\usepackage{hyperref}
\hypersetup{
    colorlinks = true,
    linkcolor  = blue!70!black,
    citecolor  = blue!70!black,
    urlcolor   = blue!70!black,
}

% ── Theorem environments ─────────────────────────────────────
\newtheorem{theorem}{Theorem}
\newtheorem{lemma}[theorem]{Lemma}
\newtheorem{proposition}[theorem]{Proposition}
\newtheorem{corollary}[theorem]{Corollary}
\theoremstyle{definition}
\newtheorem{definition}{Definition}
\newtheorem{problem}{Problem}
\theoremstyle{remark}
\newtheorem{remark}{Remark}

% ── Qubit mapping / routing macros ───────────────────────────
\newcommand{\G}{\mathcal{G}}            % coupling graph
\newcommand{\A}{\mathcal{A}}            % architecture graph
\newcommand{\C}{\mathcal{C}}            % circuit
\newcommand{\Q}{\mathcal{Q}}            % qubit set
\newcommand{\M}{\mathcal{M}}            % mapping
\newcommand{\Lay}{\mathcal{L}}          % circuit layering (ordered sequence of layers)
\newcommand{\Rout}{\mathcal{R}}         % a routing (ordered sequence of routing steps)
\newcommand{\Rmin}{R_{\min}}            % min routing steps (= circuit depth, a lower bound)
\newcommand{\SWAP}{\mathrm{SWAP}}
\newcommand{\CNOT}{\mathrm{CNOT}}
\newcommand{\depth}{\mathrm{depth}}
\newcommand{\dist}{\mathrm{dist}}       % graph distance
\newcommand{\NP}{\mathsf{NP}}
\newcommand{\todo}[1]{\textcolor{red}{\textbf{[TODO: #1]}}}

% Placeholder box for figures whose final artwork is dropped in later.
\newcommand{\figslot}[2][0.95\linewidth]{%
    \setlength{\fboxsep}{0pt}%
    \fbox{\parbox[c][#2][c]{#1}{\centering\footnotesize\itshape
        \textcolor{gray}{[ figure placeholder --- \detokenize{#1} wide ]}}}%
}

% Fixed-width variant label so the "/" lines up across regime rows.
\newcommand{\regime}[2]{\makebox[\widthof{coarse}][l]{#1} / #2}

% ── TikZ palette and tile styles (used by the overview figure) ──
\definecolor{ftblue}{HTML}{1F4E79}
\definecolor{ftorange}{HTML}{E8820C}
\definecolor{ftgreen}{HTML}{2E8B57}
\definecolor{ftgray}{HTML}{555555}
\tikzset{
  tile/.style    ={draw=ftgray!55, minimum size=5mm, inner sep=0pt},
  data/.style    ={tile, fill=ftblue!25, draw=ftblue, rounded corners=0.6pt},
  magic/.style   ={tile, fill=ftorange!85, draw=ftorange!55!black},
  free/.style    ={tile, fill=white},
  corridor/.style={tile, fill=ftgreen!35, draw=ftgreen!55!black},
  mstar/.style   ={star, star points=5, star point ratio=2.3, fill=white,
                   draw=ftorange!50!black, minimum size=2.6mm, inner sep=0pt,
                   line width=0.3pt},
  qlabel/.style  ={font=\scriptsize, text=ftblue!75!black},
  pin/.style     ={-{Stealth[length=1.6mm]}, thick, shorten >=2.5mm, shorten <=2.5mm},
}
% W x H lattice of free tiles at integer coordinates
\newcommand{\latgrid}[2]{%
  \foreach \xx in {0,...,#1}\foreach \yy in {0,...,#2}{%
     \node[free] at (\xx,\yy) {};}}
.
% ============================================================

% ── Core packages ────────────────────────────────────────────
\usepackage[T1]{fontenc}
\usepackage[utf8]{inputenc}
\usepackage{microtype}          % better line breaking
\usepackage{cite}               % IEEE-style compressed citations [1]-[3]
\usepackage{amsmath,amssymb,amsthm}
\usepackage{graphicx}
\usepackage{booktabs}           % \toprule / \midrule / \bottomrule
\usepackage{calc}               % \widthof for fixed-width regime labels
\usepackage{algorithm}
\usepackage{algpseudocode}      % algorithmicx pseudocode
\usepackage{xcolor}
\usepackage{tikz}
\usetikzlibrary{positioning,arrows.meta,calc,shapes.geometric,fit,backgrounds}
\usepackage{hyperref}
\hypersetup{
    colorlinks = true,
    linkcolor  = blue!70!black,
    citecolor  = blue!70!black,
    urlcolor   = blue!70!black,
}

% ── Theorem environments ─────────────────────────────────────
\newtheorem{theorem}{Theorem}
\newtheorem{lemma}[theorem]{Lemma}
\newtheorem{proposition}[theorem]{Proposition}
\newtheorem{corollary}[theorem]{Corollary}
\theoremstyle{definition}
\newtheorem{definition}{Definition}
\newtheorem{problem}{Problem}
\theoremstyle{remark}
\newtheorem{remark}{Remark}

% ── Qubit mapping / routing macros ───────────────────────────
\newcommand{\G}{\mathcal{G}}            % coupling graph
\newcommand{\A}{\mathcal{A}}            % architecture graph
\newcommand{\C}{\mathcal{C}}            % circuit
\newcommand{\Q}{\mathcal{Q}}            % qubit set
\newcommand{\M}{\mathcal{M}}            % mapping
\newcommand{\Lay}{\mathcal{L}}          % circuit layering (ordered sequence of layers)
\newcommand{\Rout}{\mathcal{R}}         % a routing (ordered sequence of routing steps)
\newcommand{\Rmin}{R_{\min}}            % min routing steps (= circuit depth, a lower bound)
\newcommand{\SWAP}{\mathrm{SWAP}}
\newcommand{\CNOT}{\mathrm{CNOT}}
\newcommand{\depth}{\mathrm{depth}}
\newcommand{\dist}{\mathrm{dist}}       % graph distance
\newcommand{\NP}{\mathsf{NP}}
\newcommand{\todo}[1]{\textcolor{red}{\textbf{[TODO: #1]}}}

% Placeholder box for figures whose final artwork is dropped in later.
\newcommand{\figslot}[2][0.95\linewidth]{%
    \setlength{\fboxsep}{0pt}%
    \fbox{\parbox[c][#2][c]{#1}{\centering\footnotesize\itshape
        \textcolor{gray}{[ figure placeholder --- \detokenize{#1} wide ]}}}%
}

% Fixed-width variant label so the "/" lines up across regime rows.
\newcommand{\regime}[2]{\makebox[\widthof{coarse}][l]{#1} / #2}

% ── TikZ palette and tile styles (used by the overview figure) ──
\definecolor{ftblue}{HTML}{1F4E79}
\definecolor{ftorange}{HTML}{E8820C}
\definecolor{ftgreen}{HTML}{2E8B57}
\definecolor{ftgray}{HTML}{555555}
\tikzset{
  tile/.style    ={draw=ftgray!55, minimum size=5mm, inner sep=0pt},
  data/.style    ={tile, fill=ftblue!25, draw=ftblue, rounded corners=0.6pt},
  magic/.style   ={tile, fill=ftorange!85, draw=ftorange!55!black},
  free/.style    ={tile, fill=white},
  corridor/.style={tile, fill=ftgreen!35, draw=ftgreen!55!black},
  mstar/.style   ={star, star points=5, star point ratio=2.3, fill=white,
                   draw=ftorange!50!black, minimum size=2.6mm, inner sep=0pt,
                   line width=0.3pt},
  qlabel/.style  ={font=\scriptsize, text=ftblue!75!black},
  pin/.style     ={-{Stealth[length=1.6mm]}, thick, shorten >=2.5mm, shorten <=2.5mm},
}
% W x H lattice of free tiles at integer coordinates
\newcommand{\latgrid}[2]{%
  \foreach \xx in {0,...,#1}\foreach \yy in {0,...,#2}{%
     \node[free] at (\xx,\yy) {};}}
.
% ============================================================

% ── Core packages ────────────────────────────────────────────
\usepackage[T1]{fontenc}
\usepackage[utf8]{inputenc}
\usepackage{microtype}          % better line breaking
\usepackage{cite}               % IEEE-style compressed citations [1]-[3]
\usepackage{amsmath,amssymb,amsthm}
\usepackage{graphicx}
\usepackage{booktabs}           % \toprule / \midrule / \bottomrule
\usepackage{calc}               % \widthof for fixed-width regime labels
\usepackage{algorithm}
\usepackage{algpseudocode}      % algorithmicx pseudocode
\usepackage{xcolor}
\usepackage{tikz}
\usetikzlibrary{positioning,arrows.meta,calc,shapes.geometric,fit,backgrounds}
\usepackage{hyperref}
\hypersetup{
    colorlinks = true,
    linkcolor  = blue!70!black,
    citecolor  = blue!70!black,
    urlcolor   = blue!70!black,
}

% ── Theorem environments ─────────────────────────────────────
\newtheorem{theorem}{Theorem}
\newtheorem{lemma}[theorem]{Lemma}
\newtheorem{proposition}[theorem]{Proposition}
\newtheorem{corollary}[theorem]{Corollary}
\theoremstyle{definition}
\newtheorem{definition}{Definition}
\newtheorem{problem}{Problem}
\theoremstyle{remark}
\newtheorem{remark}{Remark}

% ── Qubit mapping / routing macros ───────────────────────────
\newcommand{\G}{\mathcal{G}}            % coupling graph
\newcommand{\A}{\mathcal{A}}            % architecture graph
\newcommand{\C}{\mathcal{C}}            % circuit
\newcommand{\Q}{\mathcal{Q}}            % qubit set
\newcommand{\M}{\mathcal{M}}            % mapping
\newcommand{\Lay}{\mathcal{L}}          % circuit layering (ordered sequence of layers)
\newcommand{\Rout}{\mathcal{R}}         % a routing (ordered sequence of routing steps)
\newcommand{\Rmin}{R_{\min}}            % min routing steps (= circuit depth, a lower bound)
\newcommand{\SWAP}{\mathrm{SWAP}}
\newcommand{\CNOT}{\mathrm{CNOT}}
\newcommand{\depth}{\mathrm{depth}}
\newcommand{\dist}{\mathrm{dist}}       % graph distance
\newcommand{\NP}{\mathsf{NP}}
\newcommand{\todo}[1]{\textcolor{red}{\textbf{[TODO: #1]}}}

% Placeholder box for figures whose final artwork is dropped in later.
\newcommand{\figslot}[2][0.95\linewidth]{%
    \setlength{\fboxsep}{0pt}%
    \fbox{\parbox[c][#2][c]{#1}{\centering\footnotesize\itshape
        \textcolor{gray}{[ figure placeholder --- \detokenize{#1} wide ]}}}%
}

% Fixed-width variant label so the "/" lines up across regime rows.
\newcommand{\regime}[2]{\makebox[\widthof{coarse}][l]{#1} / #2}

% ── TikZ palette and tile styles (used by the overview figure) ──
\definecolor{ftblue}{HTML}{1F4E79}
\definecolor{ftorange}{HTML}{E8820C}
\definecolor{ftgreen}{HTML}{2E8B57}
\definecolor{ftgray}{HTML}{555555}
\tikzset{
  tile/.style    ={draw=ftgray!55, minimum size=5mm, inner sep=0pt},
  data/.style    ={tile, fill=ftblue!25, draw=ftblue, rounded corners=0.6pt},
  magic/.style   ={tile, fill=ftorange!85, draw=ftorange!55!black},
  free/.style    ={tile, fill=white},
  corridor/.style={tile, fill=ftgreen!35, draw=ftgreen!55!black},
  mstar/.style   ={star, star points=5, star point ratio=2.3, fill=white,
                   draw=ftorange!50!black, minimum size=2.6mm, inner sep=0pt,
                   line width=0.3pt},
  qlabel/.style  ={font=\scriptsize, text=ftblue!75!black},
  pin/.style     ={-{Stealth[length=1.6mm]}, thick, shorten >=2.5mm, shorten <=2.5mm},
}
% W x H lattice of free tiles at integer coordinates
\newcommand{\latgrid}[2]{%
  \foreach \xx in {0,...,#1}\foreach \yy in {0,...,#2}{%
     \node[free] at (\xx,\yy) {};}}
.
% ============================================================

% ── Core packages ────────────────────────────────────────────
\usepackage[T1]{fontenc}
\usepackage[utf8]{inputenc}
\usepackage{microtype}          % better line breaking
\usepackage{cite}               % IEEE-style compressed citations [1]-[3]
\usepackage{amsmath,amssymb,amsthm}
\usepackage{graphicx}
\usepackage{booktabs}           % \toprule / \midrule / \bottomrule
\usepackage{calc}               % \widthof for fixed-width regime labels
\usepackage{algorithm}
\usepackage{algpseudocode}      % algorithmicx pseudocode
\usepackage{xcolor}
\usepackage{tikz}
\usetikzlibrary{positioning,arrows.meta,calc,shapes.geometric,fit,backgrounds}
\usepackage{hyperref}
\hypersetup{
    colorlinks = true,
    linkcolor  = blue!70!black,
    citecolor  = blue!70!black,
    urlcolor   = blue!70!black,
}

% ── Theorem environments ─────────────────────────────────────
\newtheorem{theorem}{Theorem}
\newtheorem{lemma}[theorem]{Lemma}
\newtheorem{proposition}[theorem]{Proposition}
\newtheorem{corollary}[theorem]{Corollary}
\theoremstyle{definition}
\newtheorem{definition}{Definition}
\newtheorem{problem}{Problem}
\theoremstyle{remark}
\newtheorem{remark}{Remark}

% ── Qubit mapping / routing macros ───────────────────────────
\newcommand{\G}{\mathcal{G}}            % coupling graph
\newcommand{\A}{\mathcal{A}}            % architecture graph
\newcommand{\C}{\mathcal{C}}            % circuit
\newcommand{\Q}{\mathcal{Q}}            % qubit set
\newcommand{\M}{\mathcal{M}}            % mapping
\newcommand{\Lay}{\mathcal{L}}          % circuit layering (ordered sequence of layers)
\newcommand{\Rout}{\mathcal{R}}         % a routing (ordered sequence of routing steps)
\newcommand{\Rmin}{R_{\min}}            % min routing steps (= circuit depth, a lower bound)
\newcommand{\SWAP}{\mathrm{SWAP}}
\newcommand{\CNOT}{\mathrm{CNOT}}
\newcommand{\depth}{\mathrm{depth}}
\newcommand{\dist}{\mathrm{dist}}       % graph distance
\newcommand{\NP}{\mathsf{NP}}
\newcommand{\todo}[1]{\textcolor{red}{\textbf{[TODO: #1]}}}

% Placeholder box for figures whose final artwork is dropped in later.
\newcommand{\figslot}[2][0.95\linewidth]{%
    \setlength{\fboxsep}{0pt}%
    \fbox{\parbox[c][#2][c]{#1}{\centering\footnotesize\itshape
        \textcolor{gray}{[ figure placeholder --- \detokenize{#1} wide ]}}}%
}

% Fixed-width variant label so the "/" lines up across regime rows.
\newcommand{\regime}[2]{\makebox[\widthof{coarse}][l]{#1} / #2}

% ── TikZ palette and tile styles (used by the overview figure) ──
\definecolor{ftblue}{HTML}{1F4E79}
\definecolor{ftorange}{HTML}{E8820C}
\definecolor{ftgreen}{HTML}{2E8B57}
\definecolor{ftgray}{HTML}{555555}
\tikzset{
  tile/.style    ={draw=ftgray!55, minimum size=5mm, inner sep=0pt},
  data/.style    ={tile, fill=ftblue!25, draw=ftblue, rounded corners=0.6pt},
  magic/.style   ={tile, fill=ftorange!85, draw=ftorange!55!black},
  free/.style    ={tile, fill=white},
  corridor/.style={tile, fill=ftgreen!35, draw=ftgreen!55!black},
  mstar/.style   ={star, star points=5, star point ratio=2.3, fill=white,
                   draw=ftorange!50!black, minimum size=2.6mm, inner sep=0pt,
                   line width=0.3pt},
  qlabel/.style  ={font=\scriptsize, text=ftblue!75!black},
  pin/.style     ={-{Stealth[length=1.6mm]}, thick, shorten >=2.5mm, shorten <=2.5mm},
}
% W x H lattice of free tiles at integer coordinates
\newcommand{\latgrid}[2]{%
  \foreach \xx in {0,...,#1}\foreach \yy in {0,...,#2}{%
     \node[free] at (\xx,\yy) {};}}
